% !TEX program = xelatex
\documentclass[11pt,a4paper]{article}

\usepackage[UTF8]{ctex}
\usepackage[margin=2.2cm]{geometry}
\usepackage{amsmath, amssymb, amsthm}
\usepackage{listings}
\usepackage{xcolor}
\usepackage{tikz}
\usepackage{booktabs}
\usepackage{multirow}
\usepackage{hyperref}
\usetikzlibrary{positioning, arrows.meta, shapes.geometric, fit, calc, decorations.pathreplacing}

\hypersetup{
  colorlinks=true,
  linkcolor=blue!60!black,
  urlcolor=teal!70!black,
}

\lstdefinelanguage{Rust}{
  morekeywords={fn,let,mut,pub,struct,impl,use,mod,self,Self,Option,Some,None,return,if,else,for,while,loop,match,break,continue,as,where,trait,type,enum,move,ref,const,static,unsafe,extern,crate,super,dyn,async,await,box,in},
  sensitive=true,
  morecomment=[l]{//},
  morecomment=[s]{/*}{*/},
  morestring=[b]",
}

\lstset{
  language=Rust,
  basicstyle=\ttfamily\footnotesize,
  keywordstyle=\color{blue!60!black}\bfseries,
  commentstyle=\color{green!40!black}\itshape,
  stringstyle=\color{red!60!black},
  numbers=left,
  numberstyle=\tiny\color{gray},
  numbersep=8pt,
  frame=single,
  framerule=0.4pt,
  rulecolor=\color{gray!40},
  breaklines=true,
  breakatwhitespace=true,
  showstringspaces=false,
  tabsize=4,
  columns=fullflexible,
  keepspaces=true,
}

\title{\textbf{halfedge\,: 示例集设计文档} \\ \large \texttt{examples/}}
\author{halfedge 库}
\date{2026-07-04}

\begin{document}
\maketitle
\tableofcontents

\section{示例总览}

\texttt{examples/} 目录下提供 14 个独立可运行的示例，按功能划分，每个示例聚焦一个模块的核心用法。
所有示例通过 \texttt{cargo run --example <name>} 运行，无需额外依赖。

\begin{center}
\renewcommand{\arraystretch}{1.18}
\begin{tabular}{@{}p{3.4cm}p{2.0cm}p{5.2cm}p{5.4cm}@{}}
  \toprule
  \textbf{示例} & \textbf{模块} & \textbf{功能} & \textbf{运行命令} \\
  \midrule
  \texttt{storage\_basic}            & \texttt{storage}    & 顶点/半边/面 CRUD 与句柄有效性 & \texttt{cargo run --example storage\_basic} \\
  \texttt{build\_mesh}               & \texttt{io}         & 从顶点+面索引构建完整网格     & \texttt{cargo run --example build\_mesh} \\
  \texttt{obj\_io}                   & \texttt{io}         & OBJ 文本/文件读写往返         & \texttt{cargo run --example obj\_io} \\
  \texttt{traversal}                 & \texttt{traversal}  & 5 迭代器 + 2 边界判定函数     & \texttt{cargo run --example traversal} \\
  \texttt{topology\_ops}             & \texttt{topology\_ops} & split/flip/collapse 三大操作 & \texttt{cargo run --example topology\_ops} \\
  \texttt{extrude\_face}             & \texttt{topology\_ops} & 单/批量三角面挤出形成棱柱    & \texttt{cargo run --example extrude\_face} \\
  \texttt{geometry\_query}           & \texttt{geometry}   & 边长/面积/法向/最小内角/顶点法向 & \texttt{cargo run --example geometry\_query} \\
  \texttt{laplacian\_smooth}         & \texttt{geometry}   & 拉普拉斯平滑（顶点级+整网级） & \texttt{cargo run --example laplacian\_smooth} \\
  \texttt{point\_triangle\_distance} & \texttt{geometry}   & 点到三角形 7-Voronoi 区域距离 & \texttt{cargo run --example point\_triangle\_distance} \\
  \texttt{loop\_subdivision}         & \texttt{subdiv}     & Loop 细分规模/度数/Euler 验证 & \texttt{cargo run --example loop\_subdivision} \\
  \texttt{property}                  & \texttt{property}   & OpenMesh 风格属性系统全演示   & \texttt{cargo run --example property} \\
  \texttt{validate}                  & \texttt{validate}   & 拓扑自检：干净通过 + 破坏检测 & \texttt{cargo run --example validate} \\
  \texttt{export\_wgpu}              & \texttt{export}     & 导出 wgpu 顶点/索引缓冲       & \texttt{cargo run --example export\_wgpu} \\
  \texttt{icosphere}                 & \texttt{test\_util} & icosphere 生成与性质验证     & \texttt{cargo run --example icosphere} \\
  \bottomrule
\end{tabular}
\end{center}

\section{示例依赖关系}

\begin{center}
\resizebox{\textwidth}{!}{%
\begin{tikzpicture}[
  node distance=0.7cm and 1.0cm,
  base/.style={draw, rounded corners=2pt, minimum width=2.8cm, minimum height=0.7cm, align=center, font=\footnotesize},
  foundation/.style={base, fill=blue!10},
  builder/.style={base, fill=green!12},
  consumer/.style={base, fill=yellow!18},
  op/.style={-Stealth, thick, draw=gray!60!black}
]
  % 基础层
  \node[foundation] (storage) {\texttt{storage\_basic}};
  \node[foundation, right=of storage] (traversal) {\texttt{traversal}};

  % 构建层
  \node[builder, below=of storage] (build) {\texttt{build\_mesh}};
  \node[builder, right=of build] (obj) {\texttt{obj\_io}};
  \node[builder, right=of obj] (ico) {\texttt{icosphere}};

  % 消费层第一行
  \node[consumer, below=of build] (topo) {\texttt{topology\_ops}};
  \node[consumer, right=of topo] (extrude) {\texttt{extrude\_face}};
  \node[consumer, right=of extrude] (geo) {\texttt{geometry\_query}};
  \node[consumer, right=of geo] (lap) {\texttt{laplacian\_smooth}};

  % 消费层第二行
  \node[consumer, below=of lap] (ptd) {\texttt{point\_triangle\_distance}};
  \node[consumer, left=of ptd] (loop) {\texttt{loop\_subdivision}};
  \node[consumer, left=of loop] (prop) {\texttt{property}};
  \node[consumer, left=of prop] (val) {\texttt{validate}};
  \node[consumer, left=of val] (exp) {\texttt{export\_wgpu}};

  % 依赖箭头
  \draw[op] (storage) -- (build);
  \draw[op] (build) -- (topo);
  \draw[op] (build) -- (extrude);
  \draw[op] (build) -- (val);
  \draw[op] (build) -- (exp);
  \draw[op] (build) -- (prop);
  \draw[op] (obj) -- (build);
  \draw[op] (ico) -- (geo);
  \draw[op] (ico) -- (lap);
  \draw[op] (ico) -- (loop);
  \draw[op] (ico) -- (exp);
  \draw[op] (traversal) -- (geo);
  \draw[op] (traversal) -- (prop);

  % 图例
  \node[font=\footnotesize, below=0.5cm of exp, align=left] (legend)
    {{\color{blue!60!black}$\blacksquare$} 基础层 \quad
     {\color{green!50!black}$\blacksquare$} 构建层 \quad
     {\color{yellow!60!black}$\blacksquare$} 消费层};
\end{tikzpicture}
}
\end{center}

\section{各示例说明}

\subsection{storage\_basic：底层存储 CRUD}

演示 \texttt{MeshStorage} 的增删改查与 slotmap 的 ABA 安全保证。
关键点：
\begin{itemize}
  \item 创建 3 顶点 + 3 半边 + 1 面，手动设置 \texttt{next/prev/face} 环；
  \item 删除顶点后句柄立即失效，\texttt{contains\_vertex} 返回 \texttt{false}；
  \item 槽位复用后旧句柄不复活（版本号机制）。
\end{itemize}

\subsection{build\_mesh：从顶点面索引构建网格}

演示 \texttt{build\_mesh\_from\_vertices\_faces}，自动构建 twin/next/prev/边界环。
关键点：
\begin{itemize}
  \item 输入 4 顶点 + 2 三角形（共享 1 边），输出 10 半边（2 内部 + 8 边界）；
  \item 用 \texttt{validate\_topology} 验证构建结果；
  \item 用 \texttt{is\_boundary\_edge} 区分内部边与边界边。
\end{itemize}

\subsection{obj\_io：OBJ 读写}

演示 \texttt{parse\_obj} / \texttt{format\_obj} / \texttt{save\_obj} / \texttt{load\_obj}。
关键点：
\begin{itemize}
  \item 解析四面体 OBJ 文本，要求面绕向一致；
  \item 序列化回 OBJ 文本，文件 IO 往返一致；
  \item 支持 \texttt{v/vt/vn} 格式与负索引。
\end{itemize}

\subsection{traversal：邻域遍历}

演示 5 个迭代器 + 2 个边界判定函数。
关键点：
\begin{itemize}
  \item \texttt{VertexRing}：顶点的 outgoing 半边环（CCW 顺序）；
  \item \texttt{VertexAdjacentVerts} / \texttt{VertexAdjacentFaces}：邻居顶点 / 邻接面；
  \item \texttt{FaceHalfEdges} / \texttt{FaceVertices}：面边界环；
  \item 迭代期不持有借用，可自由 \texttt{\&mut mesh}。
\end{itemize}

\subsection{topology\_ops：三大拓扑操作}

演示 \texttt{split\_edge} / \texttt{flip\_edge} / \texttt{collapse\_edge}。
关键点：
\begin{itemize}
  \item split：顶点 $+1$，面 $+2$（内部边）或 $+1$（边界边）；
  \item flip：顶点/面/半边数不变，仅重连 4 条半边；
  \item collapse：顶点 $-1$，面 $-2$，半边 $-6$；
  \item 翻转边界边返回 \texttt{FlipOnBoundaryEdge} 错误。
\end{itemize}

\subsection{extrude\_face：三角面挤出}

演示 \texttt{extrude\_face} / \texttt{extrude\_faces} 将三角面沿给定方向
\texttt{offset} 挤出形成闭合三棱柱。
关键点：
\begin{itemize}
  \item 单面挤出：新增 3 顶点 + 7 面（1 反向顶面 + 6 侧面三角形），Euler $=2$；
  \item 不同 \texttt{offset} 方向（沿法向、斜向、反向）都能正确闭合；
  \item 零向量 \texttt{offset} 返回 \texttt{DegenerateTriangle}；
  \item 与三角形共面的 \texttt{offset}（侧面面积 $< 10^{-12}$）同样报退化；
  \item 批量挤出 \texttt{extrude\_faces} 对不相交三角形集合一次性挤出；
  \item 挤出后所有面边界环长度均为 3，\texttt{validate\_topology} 通过。
\end{itemize}

\subsection{geometry\_query：几何查询}

演示 5 个查询函数：边长、面积、法向、最小内角、顶点法向。
关键点：
\begin{itemize}
  \item icosphere(0) 所有面为等边三角形，最小内角 $= 60°$；
  \item 总面积 $\approx 4\pi$（单位球面面积）；
  \item 顶点法向与顶点位置平行（球面性质）。
\end{itemize}

\subsection{laplacian\_smooth：拉普拉斯平滑}

演示 \texttt{laplacian\_smooth\_vertex} 与 \texttt{laplacian\_smooth\_mesh}。
关键点：
\begin{itemize}
  \item 单顶点：新位置 = 邻居位置算术平均；
  \item 整网：$\lambda=0.5$，5 次迭代，重心偏移 $< 10^{-15}$（保留重心）；
  \item 平滑使球面收缩（顶点偏离单位球面）；
  \item $\lambda=0$ 时网格不变。
\end{itemize}

\subsection{point\_triangle\_distance：点到三角形距离}

演示 Ericson 7-Voronoi 区域算法。
关键点：
\begin{itemize}
  \item 面域：投影在三角形内部，距离 = 垂直高度；
  \item 顶点域 $V_A/V_B/V_C$：距离到对应顶点；
  \item 边域 $E_{AB}/E_{AC}/E_{BC}$：距离到边投影；
  \item 退化三角形（共线）退化为三顶点最近距离。
\end{itemize}

\subsection{loop\_subdivision：Loop 细分}

演示 \texttt{subdiv::loop\_subdivide} 在 icosphere 上的光滑细分效果。
关键点：
\begin{itemize}
  \item 单次细分规模：$V' = V + E$，$F' = 4F$（icosphere(1) $42 \to 162$ 顶点，$80 \to 320$ 面）；
  \item 连续多次细分保持 \texttt{check\_topology} 通过；
  \item 度数分布：原始 12 顶点保持度数 $=5$，新增边中点度数 $=6$（正则性质）；
  \item Euler 示性数 $V - E + F = 2$ 在每次细分后保持；
  \item Loop 细分是逼近型细分，单位球面偏差随细分次数减小；
  \item \texttt{validate\_topology} 完整校验无错误。
\end{itemize}

\subsection{property：OpenMesh 风格属性系统}

演示 \texttt{property::MeshProperties} 与 \texttt{PropertyHandle<T>} 的核心用法。
关键点：
\begin{itemize}
  \item 注册 \texttt{f64} 顶点权重属性，\texttt{set/get\_vertex\_prop} 读写；
  \item 类型擦除：同一容器并存 \texttt{f64 / i32 / String / bool} 多种属性；
  \item 加权拉普拉斯平滑：北极顶点权重 $5.0$，其余 $1.0$，验证等权与加权位移；
  \item 半边属性：缓存所有半边边长，统计 \texttt{min/max/avg}；
  \item 面属性：缓存面法向，与面心方向一致；
  \item \texttt{remove\_vertex\_with\_props} 删除顶点时同步清理属性；
  \item newtype 模式：同底层 \texttt{f64} 通过 \texttt{Weight}/\texttt{Temperature}
        注册两个互不冲突的属性。
\end{itemize}

\subsection{validate：拓扑自检}

演示 \texttt{validate\_topology} 的完整校验。
关键点：
\begin{itemize}
  \item 干净 icosphere(1) 通过校验（0 错误）；
  \item 人为破坏 twin 互指 $\to$ 检测 \texttt{TwinMismatch} + \texttt{SelfLoopHalfEdge}；
  \item 悬空顶点引用 $\to$ \texttt{HalfEdgeDanglingVertex}；
  \item 共线三角形 $\to$ \texttt{DegenerateFace}；
  \item 入口 origin 不匹配 $\to$ \texttt{VertexHalfEdgeInconsistent}。
\end{itemize}

\subsection{export\_wgpu：wgpu 缓冲导出}

演示 \texttt{mesh\_to\_vertex\_index\_buffers}。
关键点：
\begin{itemize}
  \item 输出 \texttt{Vec<[f32; 3]>} 顶点 + \texttt{Vec<u32>} 索引；
  \item icosphere(1)：42 顶点（504 字节）+ 240 索引（960 字节）；
  \item 索引范围合法（max index $<$ vertex count）；
  \item CCW 朝向：法向·重心 $> 0$（朝外）。
\end{itemize}

\subsection{icosphere：球体生成与性质验证}

演示 \texttt{build\_icosphere} 在不同细分级别的网格性质。
关键点：
\begin{itemize}
  \item 规模公式：$V_n = 10 \cdot 4^n + 2$，$F_n = 20 \cdot 4^n$，$E_n = 30 \cdot 4^n$；
  \item 所有顶点在单位球面上（偏差 $< 10^{-15}$）；
  \item 所有面法向朝外；
  \item icosahedron(0) 的 12 个顶点度数均为 5；
  \item Euler 公式 $V - E + F = 2$ 成立。
\end{itemize}

\section{运行方式}

\subsection{运行单个示例}

\begin{lstlisting}
cargo run --example storage_basic
cargo run --example icosphere
\end{lstlisting}

\subsection{编译所有示例（不运行）}

\begin{lstlisting}
cargo build --examples
\end{lstlisting}

\subsection{列出所有可用示例}

\begin{lstlisting}
cargo run --example  # 按 Tab 自动补全
\end{lstlisting}

\section{示例设计原则}

\begin{enumerate}
  \item \textbf{自包含}：每个示例独立可运行，不依赖其他示例；
  \item \textbf{聚焦单一功能}：每个示例只演示一个模块的核心 API；
  \item \textbf{可观测输出}：所有示例用 \texttt{println!} 打印关键数据，便于直观验证；
  \item \textbf{使用真实网格}：大多数示例以 \texttt{build\_icosphere} 作为测试网格，
        保证拓扑合法、几何有意义；
  \item \textbf{包含错误处理}：\texttt{topology\_ops} 与 \texttt{validate} 示例
        演示了错误路径（翻转边界边、检测破坏）。
\end{enumerate}

\section{典型输出示例}

\subsection{icosphere 输出片段}

\begin{lstlisting}
=== icosphere 生成器示例 ===

[细分级别与网格规模]
     n       顶点        面         半边       公式验证
     0       12       20         60          OK
     1       42       80        240          OK
     2      162      320        960          OK
     3      642     1280       3840          OK
     4     2562     5120      15360          OK

[性质 1: 单位球面] icosphere(2) 顶点距单位球面最大偏差 = 2.22e-16
[性质 5: Euler] n=0: V-E+F = 12-30+20 = 2（应为 2）
\end{lstlisting}

\subsection{topology\_ops 输出片段}

\begin{lstlisting}
[split_edge] 在中点插入新顶点 VertexId(13v1)
  顶点 12 → 13
  面 20 → 22
  validate_mesh: Ok(())
  check_topology: OK

[collapse_edge] 折叠边，合并为新顶点 VertexId(13v1)
  顶点 12 → 11
  面 20 → 18
  半边 60 → 54

错误处理：翻转边界边 → Some("禁止翻转边界边 HalfEdgeId(1v1)")
\end{lstlisting}

\subsection{loop\_subdivision 输出片段}

\begin{lstlisting}
[1] 单次 Loop 细分规模变化
  细分前 icosphere(1): V=   42 F=   80 E=  120 HE=  240
  细分后              : V=  162 F=  320 E=  480 HE=  960
  规模验证：V'=V+E (42+120=162)，F'=4F (4*80=320)
  拓扑校验：OK

[4] Euler 示性数 V - E + F 保持不变
  细分 0x: V-E+F = 12-30+20 = 2（应为 2）
  细分 1x: V-E+F = 42-120+80 = 2（应为 2）
  细分 2x: V-E+F = 162-480+320 = 2（应为 2）
  细分 3x: V-E+F = 642-1920+1280 = 2（应为 2）
\end{lstlisting}

\subsection{extrude\_face 输出片段}

\begin{lstlisting}
[1] 单个三角形挤出
  挤出前: V=  3 F=  1 E=  3 HE=  6
  挤出后: V=  6 F=  8 E= 12 HE= 24
  新增面数 = 7（应为 7：1 顶面 + 6 侧面）
  拓扑校验：OK
  Euler 示性数 V-E+F = 2（闭合三棱柱应为 2）

[3] 边界情况：零向量 / 退化 offset
  零向量 offset：DegenerateTriangle（预期行为）
  x 轴 offset：DegenerateTriangle（预期行为：侧面退化）
\end{lstlisting}

\subsection{property 输出片段}

\begin{lstlisting}
====== 1. 基础：注册并读写顶点属性 ======
已注册顶点属性 f64，已设置 42 个顶点的权重
vertex_prop_type_count = 1

====== 2. 类型擦除：同一容器并存多种类型 ======
已追加 i32 / String / bool 三种顶点属性
vertex_prop_type_count = 4

====== 4. 半边属性：缓存边长 ======
已为 240 条半边缓存长度，halfedge_prop_type_count = 1
边长统计：min=0.6071, max=0.7174, avg=0.6760 (共 240 条半边)

====== 6. 删除元素时同步清理属性 ======
删除前：
  vVertexId(1v1) exists=true, weight=Some(1.0)
删除后：
  vVertexId(1v1) exists=false, weight=None
\end{lstlisting}

\end{document}
